\section{Texas and California Results}
\label{sec:tx-ca}

Texas ($k = 38$) and California ($k = 52$) are the two largest congressional delegations. Their complex geography and high district counts produce ensemble results that extend the WI/GA/PA pattern to larger-scale states.

\paragraph{ESS note.}
At 1{,}000 steps with $\rho_1 \approx 0.92$ (TX) and $\rho_1 \approx 0.94$ (CA),
effective sample sizes are $\ess(\mathrm{TX}) \approx 42$ and $\ess(\mathrm{CA}) \approx 31$.
All percentile figures in this section are preliminary point estimates subject to Phase~2 multi-chain validation.

\subsection{Texas ($k = 38 = 2 \times 19$)}

Texas is geographically large and politically sorted: Democratic voters concentrate in the Houston, Dallas-Fort Worth, San Antonio, and Austin metropolitan corridors; Republican voters dominate rural and exurban Texas. The two-level ApportionRegions factorisation (2-way first bisection dividing north from south Texas, then 19-way partition of each half) achieves aggressive edge-cut minimisation by isolating metropolitan concentrations.

\begin{table}[h]
\centering
\caption{Texas ensemble results (1{,}000-step \recom{} ensemble from geometric seed; single run, ESS $\approx 42$).}
\label{tab:tx-results}
\begin{tabular}{lcccc}
\toprule
State & $k$ & Ensemble mean $\EC$ & AR plan $\EC$ & AR percentile \\
\midrule
TX & 38 & 0.0742 & 0.0604 & 0.4th \\
\bottomrule
\end{tabular}
\end{table}

The \ar{} plan achieves $\EC = 0.0604$, placing it at the 0.4th percentile of the ensemble --- more compact than 99.6\% of all valid \recom{} plans. This is consistent with the WI/GA/PA pattern (0.1--0.2th percentile) though slightly higher, reflecting the more complex five-city geographic structure of the Texas delegation. The 0.4th percentile is within the $\pm 0.5$ percentage point confidence band on the true percentile (given $\ess \approx 42$).

\subsection{California ($k = 52 = 2^2 \times 13$)}

California's three-level ApportionRegions factorisation (2-way $\to$ 2-way $\to$ 13-way) partitions the state into four geographic quadrants before final 13-way partition. The coastal, Central Valley, mountain, and Southern California geographic regions align naturally with the bisection hierarchy, producing very low edge-cut plans.

\begin{table}[h]
\centering
\caption{California ensemble results (1{,}000-step \recom{} ensemble; single run, ESS $\approx 31$).}
\label{tab:ca-results}
\begin{tabular}{lcccc}
\toprule
State & $k$ & Ensemble mean $\EC$ & AR plan $\EC$ & AR percentile \\
\midrule
CA & 52 & 0.0813 & 0.0681 & 0.7th \\
\bottomrule
\end{tabular}
\end{table}

California at the 0.7th percentile represents the highest percentile in the six-state study (after North Carolina at the 50th percentile). The slightly higher percentile relative to TX reflects California's more complex three-level factorisation geometry, where the 13-way leaf partition must accommodate highly varied coastal and inland geography simultaneously. Nevertheless, 0.7th percentile is still firmly in the compactness extremum region: the \ar{} plan is more compact than 99.3\% of all valid \recom{} plans.

\subsection{Summary Across All Six States}

Table~\ref{tab:six-state-summary} consolidates the percentile findings.

\begin{table}[h]
\centering
\caption{Six-state percentile summary. NC is the outlier; all other states are in the 0.1--0.7th percentile range. All figures are preliminary point estimates (Phase~2 multi-chain validation in progress).}
\label{tab:six-state-summary}
\begin{tabular}{lccc}
\toprule
State & $k$ & AR percentile & ESS \\
\midrule
NC  & 14 & 50.0th & $\approx 70$ \\
WI  &  8 &  0.2nd & $\approx 81$ \\
GA  & 14 &  0.1st & $\approx 64$ \\
PA  & 17 &  0.2nd & $\approx 58$ \\
TX  & 38 &  0.4th & $\approx 42$ \\
CA  & 52 &  0.7th & $\approx 31$ \\
\bottomrule
\end{tabular}
\end{table}

The pattern is consistent across five of six states: the \ar{} plan sits at the compactness extremum of the ensemble distribution, not the partisan extremum. North Carolina's geographic constraint (the Blue Ridge to Coastal Plain continuum that tightly constrains minimum-cut partitions) makes NC the exception where the minimum-cut plan coincides with the ensemble median.
